% Generated draft: requires mathematical and source review.
\documentclass[11pt]{article}
\usepackage[margin=25mm]{geometry}
\usepackage{fontspec}
\setmainfont{Latin Modern Roman}
\usepackage{amsmath,amssymb,mathtools,unicode-math}
\setmathfont{Latin Modern Math}
\usepackage{enumitem,needspace,xurl,hyperref}
\setlength{\parindent}{0pt}
\setlength{\parskip}{5pt}
\setlength{\emergencystretch}{3em}
\newcommand{\N}{\mathbb{N}}
\newcommand{\Z}{\mathbb{Z}}
\newcommand{\Q}{\mathbb{Q}}
\newcommand{\R}{\mathbb{R}}
\newcommand{\C}{\mathbb{C}}
\newcommand{\F}{\mathbb{F}}
\newcommand{\A}{\mathbb{A}}
\newcommand{\PP}{\mathbb{P}}
\newcommand{\E}{\mathbb{E}}
\newcommand{\eps}{\varepsilon}
\newcommand{\dd}{\,\mathrm{d}}
\begin{document}
\section*{UnsolvedMath ID: 701655}
\textbf{Source problem code:} catalogue entry
\subsection{Definitions and mathematical statement}
Use the usual convention that locally compact groups are Hausdorff, with no countability assumption. All subgroups carry the subspace topology, and quotients carry the quotient topology. A locally compact group \(G\) is amenable if, for a left Haar measure \(\mu\), there exists a positive linear functional \(m:L^\infty(G,\mu;\R)\to\R\) such that \(m(1)=1\) and \(m(L_gf)=m(f)\) for every \(g\in G\) and \(f\in L^\infty(G,\mu;\R)\), where \((L_gf)(x)=f(g^{-1}x)\).

Let \(\mathscr{AE}\) be the smallest class of locally compact groups, understood up to topological group isomorphism, satisfying the following conditions:
\begin{enumerate}
\item Every compact group and every discrete amenable group belongs to \(\mathscr{AE}\).
\item If \(G\in\mathscr{AE}\) and \(H\leq G\) is closed, then \(H\in\mathscr{AE}\).
\item If \(G\in\mathscr{AE}\) and \(N\) is a closed normal subgroup of \(G\), then the Hausdorff quotient \(G/N\) belongs to \(\mathscr{AE}\).
\item If \(G\) is a locally compact group and \((G_i)_{i\in I}\) is a family of open subgroups of \(G\), each belonging to \(\mathscr{AE}\), such that \(G=\bigcup_{i\in I}G_i\) and for every \(i,j\in I\) there exists \(k\in I\) with \(G_i\cup G_j\subseteq G_k\), then \(G\in\mathscr{AE}\).
\item If \(G\) is a locally compact group with a closed normal subgroup \(N\) such that both \(N\) and \(G/N\) belong to \(\mathscr{AE}\), then \(G\in\mathscr{AE}\).
\end{enumerate}
\textbf{Question.} Is it true that, for every locally compact group \(G\),
\[
G\text{ amenable}\quad\Longrightarrow\quad G\in\mathscr{AE}?
\]
The supplied source additionally states that a positive answer to Kourovka Notebook Problem 19.21 would imply a negative answer to this question. The statement of Problem 19.21 is not supplied, so this cross-reference is retained as a source qualification without further formulation or verification.
\subsection{Short English statement}
Must every amenable locally compact group belong to the smallest class containing all compact groups and discrete amenable groups and closed under closed subgroups, Hausdorff quotients, directed unions of open subgroups, and topological group extensions? The supplied registry labels this problem open.
\subsection{Sources}
\begin{itemize}
\item Supplied local UnsolvedMath registry record: permanent ID \(701655\), exact source problem code \texttt{catalogue entry}, title \emph{Kourovka Notebook Problem 19.107}; upstream revision \texttt{d371ee917ac9647fbb06cd6fc529deec6f3dda82}. The record was updated on \texttt{2026-10-16T14:00:00Z} and labels the problem open; this is supplied registry metadata, not an independently verified status assessment.
\item Attributed to P. Wesolek, proposed in 2018. Supplied reference: \emph{The Kourovka Notebook (Unsolved Problems in Group Theory)}, 19th issue (2018), Problem 19.107; also identified in the record with the 21st issue edition (2026), \texttt{arXiv:1401.0300v48}. The supplied collection description names E. I. Khukhro and V. D. Mazurov as editors. These bibliographic details have not been independently verified.
\item Supplied source remark: a positive answer to Problem 19.21 would yield a negative answer to Problem 19.107. No statement of Problem 19.21 is included in the supplied evidence.
\end{itemize}
\end{document}
